Разработанное приложение, основанное на WebGPU, также требует определенных шагов по настройке и инициализации. В этой секции описывается процесс создания веб-приложения с использованием TypeScript и Next.js, а также интеграция WebGPU для выполнения вычислений на GPU.

\subsection{Создание веб-приложения}

Для удобства разработки и быстроты разработки веб-приложения был выбран фреймворк Next.js, который предоставляет мощные инструменты для создания приложений на TypeScript.

В начале необходимо настроить проект Next.js с поддержкой TypeScript. Это включает в себя создание нового проекта, установку необходимых зависимостей и настройку конфигурационных файлов, линтера и прочих инструментов разработки рекоммендованных в официальной документации Next.js \cite{nextjsdocs}.

\newpage

\begin{lstlisting}[language=bash, caption=Создание нового проекта Next.js с TypeScript]
pnpm create next-app@latest my-app --yes
cd my-app
pnpm dev
\end{lstlisting}

После создания проекта, необходимо установить типы для WebGPU, чтобы обеспечить поддержку TypeScript при работе с API WebGPU. Кроме того, для создания пользовательского интерфейса с динамическими параметрами симуляции будет использоваться библиотека \texttt{lil-gui}.

\begin{lstlisting}[language=bash, caption=Установка зависимостей]
pnpm add -D @webgpu/types # Установка типов для WebGPU
pnpm add lil-gui          # Установка библиотеки lil-gui
\end{lstlisting}

\subsection{Подключение WebGPU}

Для веб-приложения были разработаны React-компоненты, один из которых является точкой входа для инициализации WebGPU и создания контекста рендеринга. В этом компоненте происходит вызов функций инициализации WebGPU, а также передача параметров мыши (координаты и скорость) для удобного взаимодействия с симуляцией.

\begin{lstlisting}[language=ES6,caption=Часть React-компонента для инициализации WebGPU]
useEffect(() => {
    const canvas = canvasRef.current
    if (!canvas) return

    setupMouseListener(canvas, params.mouse)

    // This will hold the cleanup function returned by main()
    let cleanup: (() => void) | undefined

    main(canvas, params)
      .then((c) => {
        // Store the cleanup function for later
        cleanup = c
      })
      .catch((err) => {
        console.error(err)
      })

    return () => {
      if (cleanup) cleanup()
    }
  }, []) // only initialize WebGPU once
}
\end{lstlisting}

Функция \texttt{main} отвечает за инициализацию WebGPU, создание устройства GPU, выделение буферов и настройку вычислительных шейдеров. Она также запускает основной цикл рендеринга и обновления симуляции.

Таким образом, веб-приложение на Next.js соединяется с WebGPU, позволяя эффективно использовать возможности GPU для выполнения вычислительных задач, связанных с симуляцией жидкости.